% Lamb_Anu.tex --- master file, complete hand calculation.
% Compile: pdflatex, bibtex, pdflatex, pdflatex.
% \showworkedtrue  -> master with the working shown
% \showworkedfalse -> scholar version, results and sources only
\documentclass[11pt]{article}
\usepackage[T1]{fontenc}
\usepackage{lmodern}
\usepackage{amsmath,amssymb}
\usepackage{amsthm}
\usepackage{booktabs}
\usepackage{array}
\usepackage[a4paper,margin=2.2cm]{geometry}
\usepackage[hidelinks]{hyperref}

\newif\ifshowworked
\showworkedtrue
\makeatletter
\newenvironment{worked}{%
  \ifshowworked\par\smallskip\noindent\textit{Working.}\par\smallskip\else
  \setbox\z@\vbox\bgroup\fi}{\ifshowworked\par\smallskip\else\egroup\fi}
\makeatother

\newcommand{\ii}{\mathrm{i}}
\newcommand{\Ncal}{\mathcal{N}}
\newcommand{\Gcal}{\mathcal{G}}
\setlength{\parskip}{3pt}
\setlength{\parindent}{0pt}

\title{\vspace{-1.2cm}Lamb waves in a magneto-electro-elastic plate:\\
pyroelectric and pyromagnetic coupling, size effect and a heat equation with one lag}
\author{Dr.\ Vipin \\ Gurugram University}
\date{30 October 2026}

\begin{document}
\maketitle

\begin{abstract}
\noindent
We set up and reduce the plane-strain Lamb problem in a plate of a magneto-electro-elastic composite of the
barium titanate and cobalt ferrite family, poled along the thickness, with free faces at $z=\pm h$. Five
fields are carried: the displacements $u_x,u_z$, the electric potential $\varphi$, the magnetic potential
$\psi$ and the temperature rise $\theta$. The mechanical stress follows the nonlocal strain gradient law on
its elastic part; the piezoelectric, piezomagnetic, dielectric, magnetic and magnetoelectric couplings are
local. The heat equation is of Moore--Gibson--Thompson type in its base form and carries the pyroelectric
and pyromagnetic work together with their electrocaloric and magnetocaloric partners. Effective constants of the composite are carried as a band from two homogenisation routes. The five field equations
are reduced to a $5\times5$ matrix in the thickness exponent, its characteristic polynomial is counted and
its parity decided, the amplitude ratios are given through the cofactors, the seven face conditions are
derived with open and closed electric and magnetic faces, and the dispersion matrix is assembled. Three
hand limits close the file.
\end{abstract}

\noindent\textbf{Keywords:} guided waves; magneto-electro-elastic plate; magnetoelectric coefficient;
pyroelectric and pyromagnetic coupling; nonlocal strain gradient elasticity.

\section{Introduction}

A magneto-electro-elastic material is one in which a strain produces both an electric polarisation and a
magnetisation, and in which an electric field can produce a magnetic response and the other way round. No
single phase does this well; the usual route is a composite of a piezoelectric phase, barium titanate in
our case, and a piezomagnetic one, cobalt ferrite, sintered together. The magnetoelectric coefficient of
the composite is a product property, present only because the two phases are in contact, and it is much
larger than the product of the two bulk effects would suggest.

Guided waves in magneto-electro-elastic plates have been computed, and the plate with a temperature field
has been computed too, though separately. Xiao, Han, Liu and Li solve guided waves in a functionally graded
magneto-electro-elastic plate with no thermal field at all \cite{Xiao2016}. Kondaiah, Shankar and Ganesan
keep both the pyroelectric and the pyromagnetic coefficients in a magneto-electro-elastic plate
\cite{Kondaiah2013}, and Tocci Monaco, Fantuzzi, Fabbrocino and Luciano keep both in a nonlocal strain
gradient nanoplate with hygro-thermal loads \cite{Tocci2021}; those two are free-vibration, buckling and
static problems, not complex-wavenumber dispersion. Matlam and Kedar put a Moore--Gibson--Thompson heat
equation with memory and nonlocality into an electro-magneto-thermoelastic medium \cite{Matlam2026}, and
Kakar and Kakar treat edge waves in an electromagnetic thermoelastic plate \cite{Kakar2015}; both were
available to us only as title and abstract.

The constitutive law of a linear magneto-electro-elastic medium is standard and is written in the layered
plate setting by Chen, Pan and Chen \cite{Chen2007}; the sign convention for the magnetoelectric
coefficient and the form of the plate conditions are taken from the way that literature writes them, and
the scholar must open the paper before quoting an equation number, because only the bibliographic data
came through for me.

What is claimed here is a plate problem in which the heat equation carries a relaxation time and both the
pyroelectric and the pyromagnetic couplings are present, with the mechanical size effect, and with the
uncertainty of the effective magnetoelectric coefficient carried openly as a band between two
homogenisation routes. The corner is narrow: the paper has to show, on one figure, that the attenuation it
predicts is larger than the spread between the two routes. If it is not, the honest outcome is a negative
result about the measurability of the effect, and that is a result.

\section{Setting and assumptions}\label{sec:setting}

The plate occupies $-h\le z\le h$, is unbounded in $x$ and $y$, and is poled along $z$, so the effective
medium is transversely isotropic with the axis along the thickness. Plane strain applies, $u_y=0$,
$\partial_y=0$. The electric and magnetic potentials give the fields by $E_i=-\varphi_{,i}$ and
$H_i=-\psi_{,i}$, and one plane wave travels along $x$,
\begin{equation}
\bigl(u_x,u_z,\varphi,\psi,\theta\bigr)
=\bigl(\hat u_x,\hat u_z,\hat\varphi,\hat\psi,\hat\theta\bigr)(z)\;
\exp\bigl[\ii(kx-\omega t)\bigr],
\qquad k=k_r+\ii k_i .
\end{equation}

The assumptions: linear magneto-electro-elasticity at the reference temperature with effective constants
of the composite, held fixed; a small temperature rise; no free charge and no free current in the bulk, so
both Gauss laws are sourceless; no pre-stress; and the two size lengths on the elastic part of the
mechanical stress only. The last point is a modelling choice with a reason: it keeps $D_i$ and $B_i$
second order in $z$, so the electric and magnetic problems remain Poisson-like and need no extra face
condition, and it keeps the size effect mechanical, so the extra conditions are mechanical too.

\section{Notation}\label{sec:notation}

\begin{center}
\small
\begin{tabular}{p{2.2cm}p{7.8cm}p{4.0cm}}
\toprule
symbol & meaning in words & unit\\
\midrule
$2h$ & thickness; faces at $z=\pm h$ & m\\
$k=k_r+\ii k_i$ & complex wavenumber; $k_i>0$ is decay & m$^{-1}$\\
$\eta$ & exponent through the thickness, one per partial wave & m$^{-1}$\\
$u_x,u_z$ & displacements & m\\
$\varphi,\psi$ & electric and magnetic potentials & V, A\\
$\theta$ & temperature rise over $T_0$ & K\\
$C_{11},C_{13},C_{33},C_{44}$ & effective elastic moduli of the composite & Pa\\
$e_{31},e_{33},e_{15}$ & effective piezoelectric constants & C\,m$^{-2}$\\
$h_{31},h_{33},h_{15}$ & effective piezomagnetic constants (written $h$ so that $q$ stays a charge) & N\,A$^{-1}$m$^{-1}$\\
$\epsilon_{11},\epsilon_{33}$ & effective dielectric permittivities & F\,m$^{-1}$\\
$\mu_{11},\mu_{33}$ & effective magnetic permeabilities & N\,A$^{-2}$\\
$\alpha_{11},\alpha_{33}$ & effective magnetoelectric coefficients, unit s\,m$^{-1}$ in SI & s\,m$^{-1}$\\
$\beta_1,\beta_3$ & thermal stress coefficients & Pa\,K$^{-1}$\\
$p_3$ & effective pyroelectric constant & C\,m$^{-2}$K$^{-1}$\\
$m_3$ & effective pyromagnetic constant & N\,A$^{-1}$m$^{-1}$K$^{-1}$\\
$v$ & volume fraction of the piezoelectric phase & --\\
$\rho,c_p,k_{\rm th}$ & density, specific heat, conductivity & kg\,m$^{-3}$, J\,kg$^{-1}$K$^{-1}$, W\,m$^{-1}$K$^{-1}$\\
$\tau_q,\tau_T$ & lag of the heat equation and its thermal companion & s\\
$l_{\rm nl},l_g$ & nonlocal and gradient lengths & m\\
\bottomrule
\end{tabular}
\end{center}

The magnetoelectric coefficient deserves its unit spelled out, because it is where sign and unit slips live. With $D_i=e_{ikl}\varepsilon_{kl}+\epsilon_{ik}E_k+\alpha_{ik}H_k+p_i\theta$, the term
$\alpha H$ must have the unit of $D$, which is C\,m$^{-2}$; $H$ is in A\,m$^{-1}$, so $\alpha$ is in
C\,A$^{-1}$m$^{-1}$, and since the ampere is a coulomb per second this is s\,m$^{-1}$. A positive
$\alpha_{33}$ means that a magnetic field along the axis adds electric displacement along the same axis.
The sign depends on which phase is taken as the matrix and on the poling direction, so both must be stated
with the number.

\section{Source equations}\label{sec:source}

Kinematics, as usual,
\begin{equation}
\varepsilon_{ij}=\tfrac12(u_{i,j}+u_{j,i}),\qquad E_i=-\varphi_{,i},\qquad H_i=-\psi_{,i}.
\end{equation}

The constitutive relations of a linear magneto-electro-elastic medium with thermal terms,
\begin{align}
\bigl(1-l_{\rm nl}^2\nabla^2\bigr)t_{ij}&=\sigma_{ij},&
\sigma_{ij}&=C_{ijkl}\bigl(1-l_g^2\nabla^2\bigr)\varepsilon_{kl}
-e_{kij}E_k-h_{kij}H_k-\beta_{ij}\theta,\label{eq:C1}\\
D_i&=e_{ikl}\varepsilon_{kl}+\epsilon_{ik}E_k+\alpha_{ik}H_k+p_i\theta,\label{eq:C2}\\
B_i&=h_{ikl}\varepsilon_{kl}+\alpha_{ik}E_k+\mu_{ik}H_k+m_i\theta .\label{eq:C3}
\end{align}
The mechanical part without the size effect and without temperature is the standard law of the layered
plate literature \cite{Chen2007}; the pyroelectric and pyromagnetic terms are kept in plate work by
\cite{Kondaiah2013,Tocci2021}, and the fact that they are kept is worth stating because an older claim
that they are usually dropped is not true. The size-effect pair is from Lim, Zhang and Reddy through the
introduction of Barretta and Marotti de Sciarra \cite[Eqs.~(18)--(19)]{Barretta2018}; the Lim paper itself
was closed to us \cite{Lim2015}.

The field equations,
\begin{align}
\bigl(1-l_{\rm nl}^2\nabla^2\bigr)\rho\,\ddot u_i&=\sigma_{ij,j},\label{eq:mom}\\
D_{i,i}&=0,\qquad B_{i,i}=0,\label{eq:gauss}\\
\Bigl(1+\tau_q\frac{\partial}{\partial t}\Bigr)
\Bigl[\rho c_p\dot\theta+T_0\beta_{ij}\dot\varepsilon_{ij}
+T_0p_i\dot E_i+T_0m_i\dot H_i\Bigr]&=k_{\rm th}\nabla^2\theta .\label{eq:heat}
\end{align}
The heat equation is the Moore--Gibson--Thompson base form with the conductivity rate put to zero, the same
operator used in this literature \cite{Matlam2026,Gupta2025}; I read the title and the abstract of both.
The electrocaloric and magnetocaloric terms are the reciprocal partners of the pyroelectric and
pyromagnetic terms: dropping them while keeping $p_i$ and $m_i$ would break reciprocity. That is a
textbook statement, not a discovery, and it belongs in one sentence of the model section.

\section{Calculation A: constitutive relations in components}\label{sec:calcA}

\begin{worked}
The effective medium is transversely isotropic with the axis along $z$, so the moduli are those of class
$6mm$: $C_{11}=C_{22}$, $C_{13}=C_{23}$, $C_{33}$, $C_{44}=C_{55}$, $C_{66}=(C_{11}-C_{12})/2$, and every
mixed element vanishes. The piezoelectric tensor has $e_{31}=e_{32}$, $e_{33}$, $e_{15}=e_{24}$; the
piezomagnetic tensor has the same shape with $h$ in place of $e$; the dielectric and magnetic tensors have
two entries each; the magnetoelectric tensor of this class is diagonal in the plane,
$\alpha_{11}=\alpha_{22}$ and $\alpha_{33}$. Plane strain removes $y$.
\end{worked}

\begin{align}
\sigma_{xx}&=\Gcal\bigl(C_{11}\varepsilon_{xx}+C_{13}\varepsilon_{zz}\bigr)-e_{31}E_z-h_{31}H_z-\beta_1\theta,\\
\sigma_{zz}&=\Gcal\bigl(C_{13}\varepsilon_{xx}+C_{33}\varepsilon_{zz}\bigr)-e_{33}E_z-h_{33}H_z-\beta_3\theta,\\
\sigma_{xz}&=2C_{44}\Gcal\,\varepsilon_{xz}-e_{15}E_x-h_{15}H_x,\\
D_x&=e_{15}\,2\varepsilon_{xz}+\epsilon_{11}E_x+\alpha_{11}H_x,&
D_z&=e_{31}\varepsilon_{xx}+e_{33}\varepsilon_{zz}+\epsilon_{33}E_z+\alpha_{33}H_z+p_3\theta,\\
B_x&=h_{15}\,2\varepsilon_{xz}+\alpha_{11}E_x+\mu_{11}H_x,&
B_z&=h_{31}\varepsilon_{xx}+h_{33}\varepsilon_{zz}+\alpha_{33}E_z+\mu_{33}H_z+m_3\theta .
\end{align}

\section{Calculations B and C: the five equations and the matrix}\label{sec:calcBC}

With $\nabla^2\to L=\eta^2-k^2$,
\begin{equation}
\Ncal(\eta)=1+l_{\rm nl}^2(k^2-\eta^2),\qquad \Gcal(\eta)=1+l_g^2(k^2-\eta^2),
\end{equation}
and amplitudes $\mathbf a=(\hat u_x,\hat u_z,\hat\varphi,\hat\psi,\hat\theta)^T$, the five equations are
$M(\eta)\mathbf a=\mathbf 0$.

\subsection*{Row 1: momentum in $x$}

\begin{worked}
With $\varepsilon_{xx}=\ii kX$, $\varepsilon_{zz}=\eta Z$, $2\varepsilon_{xz}=\eta X+\ii kZ$,
$E_x=-\ii k\Phi$, $E_z=-\eta\Phi$, $H_x=-\ii k\Psi$, $H_z=-\eta\Psi$:
$\partial_x\sigma_{xx}$ gives $-\Gcal C_{11}k^2X$ from $u_x$, $\ii k\eta\Gcal C_{13}Z$ from $u_z$,
$\ii k\eta e_{31}\Phi$ and $\ii k\eta h_{31}\Psi$ from the two potentials (signs from $-e_{31}E_z$ with
$E_z=-\eta\Phi$), and $-\ii k\beta_1\Theta$ from the temperature.
$\partial_z\sigma_{xz}$ gives $\Gcal C_{44}\eta^2X$, $\ii k\eta\Gcal C_{44}Z$, $\ii k\eta e_{15}\Phi$,
$\ii k\eta h_{15}\Psi$.
Add, and move $-\Ncal\rho\omega^2X$ across.
\end{worked}

\begin{equation}
M_{11}=-\Ncal\rho\omega^2+\Gcal(C_{11}k^2-C_{44}\eta^2),\quad
M_{12}=-\ii k\eta\Gcal(C_{13}+C_{44}),\quad
M_{13}=-\ii k\eta(e_{31}+e_{15}),\quad
M_{14}=-\ii k\eta(h_{31}+h_{15}),\quad
M_{15}=\ii k\beta_1 .
\end{equation}

\subsection*{Row 2: momentum in $z$}

\begin{worked}
$\partial_x\sigma_{xz}$: $\ii k\eta\Gcal C_{44}X$, $-\Gcal C_{44}k^2Z$, $-e_{15}k^2\Phi$, $-h_{15}k^2\Psi$.
$\partial_z\sigma_{zz}$: $\ii k\eta\Gcal C_{13}X$, $\Gcal C_{33}\eta^2Z$, $e_{33}\eta^2\Phi$,
$h_{33}\eta^2\Psi$, $-\eta\beta_3\Theta$. Collect.
\end{worked}

\begin{equation}
M_{21}=-\ii k\eta\Gcal(C_{13}+C_{44}),\quad
M_{22}=-\Ncal\rho\omega^2+\Gcal(C_{44}k^2-C_{33}\eta^2),\quad
M_{23}=e_{15}k^2-e_{33}\eta^2,\quad
M_{24}=h_{15}k^2-h_{33}\eta^2,\quad
M_{25}=\eta\beta_3 .
\end{equation}

\subsection*{Row 3: electric Gauss law}

\begin{worked}
$D_x=e_{15}(\eta X+\ii kZ)-\ii k\epsilon_{11}\Phi-\ii k\alpha_{11}\Psi$;
$D_z=\ii ke_{31}X+\eta e_{33}Z-\eta\epsilon_{33}\Phi-\eta\alpha_{33}\Psi+p_3\Theta$.
$\partial_xD_x$: $\ii k\eta e_{15}X$, $-e_{15}k^2Z$, $-\epsilon_{11}(-k^2)\Phi\to k^2\epsilon_{11}\Phi$,
$k^2\alpha_{11}\Psi$.
$\partial_zD_z$: $\ii k\eta e_{31}X$, $\eta^2e_{33}Z$, $-\eta^2\epsilon_{33}\Phi$, $-\eta^2\alpha_{33}\Psi$,
$\eta p_3\Theta$.
Add and set to zero. There is no free charge.
\end{worked}

\begin{equation}
M_{31}=\ii k\eta(e_{15}+e_{31}),\quad
M_{32}=e_{33}\eta^2-e_{15}k^2,\quad
M_{33}=\epsilon_{33}\eta^2-\epsilon_{11}k^2,\quad
M_{34}=\alpha_{33}\eta^2-\alpha_{11}k^2,\quad
M_{35}=\eta p_3 .
\end{equation}

\subsection*{Row 4: magnetic Gauss law}

\begin{worked}
The same steps with $D\to B$, $e\to h$, $\epsilon\to\mu$, $p_3\to m_3$, and $\alpha$ in the same place;
the cross entry is the same $\alpha$ because the magnetoelectric tensor is symmetric in the energy, which
is the Onsager symmetry for this coupling and is why one symbol serves both rows.
\end{worked}

\begin{equation}
M_{41}=\ii k\eta(h_{15}+h_{31}),\quad
M_{42}=h_{33}\eta^2-h_{15}k^2,\quad
M_{43}=\alpha_{33}\eta^2-\alpha_{11}k^2,\quad
M_{44}=\mu_{33}\eta^2-\mu_{11}k^2,\quad
M_{45}=\eta m_3 .
\end{equation}

\subsection*{Row 5: the heat equation}

\begin{worked}
Put $f=1-\ii\omega\tau_q$. Bracket: $-\ii\omega\rho c_p\Theta$; the thermoelastic work gives
$T_0\omega k\beta_1X-\ii\omega T_0\beta_3\eta Z$; the electrocaloric term $T_0p_3\dot E_z=
\ii\omega T_0p_3\eta\Phi$; the magnetocaloric term $T_0m_3\dot H_z=\ii\omega T_0m_3\eta\Psi$.
Multiply by $f$, subtract $k_{\rm th}L\Theta$.
\end{worked}

\begin{equation}
M_{51}=fT_0\omega k\beta_1,\quad
M_{52}=-\ii\omega fT_0\beta_3\eta,\quad
M_{53}=\ii\omega fT_0p_3\eta,\quad
M_{54}=\ii\omega fT_0m_3\eta,\quad
M_{55}=-\ii\omega f\rho c_p-k_{\rm th}L .
\end{equation}

\section{Calculation D: degree and parity}\label{sec:calcD}

\begin{worked}
Degrees by entry: row 1 is four, three, three, three, one; row 2 is three, four, two, two, one; row 3 is
one, two, two, two, one; row 4 is one, two, two, two, one; row 5 is zero, one, one, one, two. The diagonal
sum is $4+4+2+2+2=14$ and no permutation improves on it, so the characteristic polynomial has degree $14$
when $l_g\ne0$ and $10$ when $l_g=0$.
For the parity, look for signs $\sigma_i,\tau_j$ with $M_{ij}(-\eta)=\sigma_i\tau_jM_{ij}(\eta)$.
Set $\tau_1=+1$: the odd entries $M_{12},M_{13},M_{14}$ give $\sigma_1\tau_2=\sigma_1\tau_3=\sigma_1\tau_4=-1$
and the even $M_{15}$ gives $\sigma_1\tau_5=+1$. Take $\sigma_1=+1$, so $\tau_2=\tau_3=\tau_4=-1$ and
$\tau_5=+1$. Row 2 then fixes $\sigma_2=-1$, and its even entries $M_{23},M_{24}$ need
$\sigma_2\tau_3=\sigma_2\tau_4=+1$, true; its odd $M_{25}$ needs $\sigma_2\tau_5=-1$, true. Row 3 fixes
$\sigma_3=-1$ from the odd $M_{31}$; the odd $M_{35}$ needs $\sigma_3\tau_5=-1$, true; the even
$M_{33},M_{34}$ are satisfied. Row 4 gives $\sigma_4=-1$ the same way. Row 5: the odd
$M_{52},M_{53},M_{54}$ need $\sigma_5\tau_j=-1$, which holds with $\sigma_5=+1$; the even $M_{51},M_{55}$
are fine. The matching closes.
\end{worked}

So the determinant is even in $\eta$ for all settings, the roots come in pairs, and the plate modes split
into a symmetric and an antisymmetric family. Each family can be posed with seven roots and seven rows, and
the two routes must agree. The magnetoelectric coupling does not break the split, because it enters as
$\alpha_{33}\eta^2$ and $\alpha_{11}k^2$, both even, in both the electric and the magnetic Gauss rows.

\section{Calculation E: amplitude ratios and the solution}\label{sec:calcE}

\begin{worked}
For each root, take a column of the cofactor matrix $C_{ij}=(-1)^{i+j}\mathcal M_{ji}$, which is a null
vector of $M(\eta_m)$ because $M\,C=(\det M)I=0$ at the root. The safest column is the one belonging to the
row with the largest norm at that root, and the arbitrary scale is absorbed in the amplitude $A_m$.
\end{worked}

\begin{equation}
\begin{pmatrix}u_x\\u_z\\\varphi\\\psi\\\theta\end{pmatrix}
=\sum_{m=1}^{14}A_m
\begin{pmatrix}X_m\\Z_m\\\Phi_m\\\Psi_m\\\Theta_m\end{pmatrix}
e^{\eta_mz}e^{\ii(kx-\omega t)} .
\end{equation}

\section{Calculation F: the seven face conditions}\label{sec:calcF}

\begin{worked}
Seven conditions per face: the normal and shear tractions with the surface-gradient term and the nonlocal
operator inverted root by root, the two double tractions, the normal electric displacement, the normal
magnetic induction, and the temperature condition. The traction and double-traction construction is the
one of \cite[Sec.~2, Eq.~(14)]{Charal2019}.
Two switches must be stated in the paper, not buried: an electrically open face uses $D_z=0$, a closed one
uses $\Phi=0$; a magnetically open face uses $B_z=0$, a closed one uses $\Psi=0$. The four combinations are
all defensible and they give different attenuation, so the paper must name the one it uses and show one
figure for the extreme pair.
\end{worked}

\begin{align}
b_1&=\frac{1}{\Ncal_m}\Bigl[\hat\sigma_{zz}-l_g^2\ii k\eta\,C_{44}(\eta X_m+\ii kZ_m)\Bigr],&
b_2&=\frac{1}{\Ncal_m}\Bigl[\hat\sigma_{xz}-l_g^2\ii k\eta\,(\ii kC_{11}X_m+\eta C_{13}Z_m)\Bigr],\\
b_3&=\frac{l_g^2}{\Ncal_m}\eta\,C_{44}(\eta X_m+\ii kZ_m),&
b_4&=\frac{l_g^2}{\Ncal_m}\eta\,(\ii kC_{13}X_m+\eta C_{33}Z_m),\\
b_5&=\ii ke_{31}X_m+\eta e_{33}Z_m-\eta\epsilon_{33}\Phi_m-\eta\alpha_{33}\Psi_m+p_3\Theta_m,&
b_6&=\ii kh_{31}X_m+\eta h_{33}Z_m-\eta\alpha_{33}\Phi_m-\eta\mu_{33}\Psi_m+m_3\Theta_m,\\
b_7&=\eta\Theta_m .&&
\end{align}

\section{Calculation G: the dispersion matrix}\label{sec:calcG}

\begin{equation}
\mathbf B(k,\omega)=
\begin{pmatrix}b_r(\eta_m)e^{\eta_mh}\\ b_r(\eta_m)e^{-\eta_mh}\end{pmatrix}_{r=1\dots7,\;m=1\dots14},
\qquad \det\mathbf B(k,\omega)=0 ,
\end{equation}
solved for complex $k$ at real $\omega$, with $k_i$ the attenuation per unit length. When $l_g=0$ the rows
$b_3,b_4$ drop and ten roots are used.

\section{Calculation H: three hand limits}\label{sec:calcH}

\subsection*{The magnetoelectric coupling off}

\begin{worked}
Set $\alpha_{11}=\alpha_{33}=0$: the electric and magnetic Gauss rows decouple. The electric row keeps
$e$ and $\epsilon$, the magnetic row keeps $h$ and $\mu$, and the mechanical row carries the sum of the two
stiffenings. That is the piezoelectric plus piezomagnetic plate, and its dispersion must be the one a
piezoelectric plate with an extra stiffening gives.
\end{worked}

\subsection*{Piezoelectric and piezomagnetic coupling both off}

\begin{worked}
With $e=h=\alpha=p=m=0$ the two Gauss rows reduce to $\epsilon_{ii}$ and $\mu_{ii}$ terms and force the
potentials to vanish; the mechanical rows reduce to the thermoelastic plate with the size effect. The
dispersion must then be the thermoelastic Rayleigh--Lamb problem.
\end{worked}

\subsection*{The thin-plate low-frequency limit and the effective magnetoelectric coefficient}

\begin{worked}
For a thin plate at low frequency the fields through the thickness are nearly uniform, so the face
conditions reduce to the requirement that the average electric displacement and the average magnetic
induction vanish when the faces are open. Eliminating the two potentials between those two conditions
leaves a coupled pair in which the strain drives an electric displacement through $e$ and the magnetic
field, and a magnetic induction through $h$ and the electric field, with the magnetoelectric terms
$\alpha$ closing the loop. Solving that $2\times2$ system gives the effective response: the ratio of the
induced magnetic induction to the applied electric field at zero strain, which is the composite
magnetoelectric coefficient of the plate. It is a function of $\alpha$, of $e$ and $h$, and of the two
permittivities; the algebra is the scholar's to complete, and the check is that it reduces to $\alpha$ when
$e=h=0$.
\end{worked}

\section{Homogenisation: two routes and the band}\label{sec:homog}

The volume fraction $v$ of the piezoelectric phase enters only through the effective constants. Two routes
are carried, and they are not equivalent, so both appear on every figure.

\emph{Route 1, the rule of mixtures.} Each constant is the volume average of the two phases,
$\bar C=vC^{p}+(1-v)C^{m}$ and the same for $\epsilon,\mu,\rho$; the coupling constants are weighted in the
same way, and the magnetoelectric coefficient is taken as
$\bar\alpha=v(1-v)\,\alpha_{0}$ with $\alpha_0$ a product of the two phase couplings, so that it vanishes
at both ends as it must.

\emph{Route 2, the Mori--Tanaka estimate.} One phase is treated as an ellipsoidal inclusion in the other,
and the effective constants follow from the strain and field concentration tensors; the magnetoelectric
coefficient is then not a simple product and is not symmetric in the choice of matrix, which is why the
choice must be stated.

Both are textbook composites estimates; neither is a fit. The spread between them is the band, and the
condition the paper has to meet is that the predicted attenuation be wider than the band. That is stated as
a test, not as a hope, and the week-nine decision rests on it.

\section{Self-test list}\label{sec:selftest}

\begin{enumerate}
\item All couplings off, both size lengths zero: the classical Rayleigh--Lamb plate, both families.
\item $\alpha=0$ with $e,h\ne0$: the decoupled piezoelectric plus piezomagnetic plate.
\item $e=h=\alpha=0$: the thermoelastic plate with the size effect.
\item Both routes at $v\to0$ and $v\to1$: the composite constants must approach the pure phases.
\item The parity statement of Calculation D, checked from the roots and from the two family matrices.
\item $l_g\to0$ and $l_{\rm nl}\to0$ separately, continuously, with the gradient length stiffening and the
nonlocal length softening.
\item Row count against root count, in both settings of $l_g$.
\item The four electric and magnetic face combinations: state which is used, and check that the closed face
gives $\Phi_m=0$ as a row.
\end{enumerate}

\section{Work plan}\label{sec:plan}

Weeks one to two: read \cite{Chen2007} and fix the sign convention of $\alpha$ and the plate conditions on
the page; read \cite{Matlam2026} for the heat equation as written and for whether the electrocaloric and
magnetocaloric terms are kept. Weeks three to four: the matrix by hand, the roots, the parity, the family
matrices. Weeks five to six: the three hand limits and the elastic check. Weeks seven to eight: the two
homogenisation routes and the band, with the material constants of both phases on the record. Weeks nine to
ten: dispersion and attenuation for the four face combinations, and the band-versus-effect figure that
decides the target journal. Weeks eleven to twelve: writing and the citation check.

\bibliographystyle{unsrt}
\bibliography{Lamb_Anu}

\end{document}
